\documentclass[12pt,a4paper,oneside]{article}             % Single-side
%\documentclass[11pt,a4paper,twoside,openright]{report}  % Duplex

%\PassOptionsToPackage{chapternumber=Huordinal}{magyar.ldf}
\usepackage{t1enc}
\usepackage[latin2]{inputenc}
\usepackage[magyar]{babel}
\usepackage[T1]{fontenc}
\usepackage[table]{xcolor}
\usepackage{makecell}
\usepackage{array}
\usepackage{mathtools}
\usepackage{amsmath}
\usepackage{amssymb}
\usepackage{enumerate}
\usepackage[thmmarks]{ntheorem}
\usepackage{graphics}
\usepackage{epsfig}
\usepackage{listings}
\usepackage{color}
%\usepackage{fancyhdr}
\usepackage{lastpage}
\usepackage{anysize}
\usepackage{sectsty}
\usepackage{setspace}  % Ettol a tablazatok, abrak, labjegyzetek maradnak 1-es sorkozzel!
\usepackage[hang]{caption}
\usepackage{hyperref}
\usepackage{indentfirst}
\usepackage[normalem]{ulem}
\usepackage{ctable}
\usepackage{cleveref}
\usepackage{mathabx}
\usepackage{verbatim}
\usepackage{tocloft}
\usepackage{float}
\usepackage{rotating}
\usepackage{pdflscape}
\usepackage{lipsum}
\usepackage{everypage}

%\usepackage[paper=portrait,pagesize]{typearea}

%--------------------------------------------------------------------------------------
% Code sytles
%--------------------------------------------------------------------------------------

% Python -----------------------------------------------------
\definecolor{weborange}{RGB}{255,165,0}
\lstdefinestyle{Python}{
    language        = Python,
    basicstyle      = \fontsize{11}{11}\ttfamily,
		%backgroundcolor=\color{white!10}
    keywordstyle    = \color{weborange},
    emphstyle				= \color{blue},
    keywordstyle    = [2] \color{teal}, % just to check that it works
    stringstyle     = \color{green},
    commentstyle    = \color{red}\ttfamily
}
\lstnewenvironment{pycode}{
\lstset{style=Python}
\lstset{literate={'}{{'}}1}}{}


% Python -----------------------------------------------------
\definecolor{mygreen}{RGB}{28,172,0} % color values Red, Green, Blue
\definecolor{mylilas}{RGB}{170,55,241}

\lstdefinestyle{Matlab}{
		language=Matlab,%
    basicstyle=\fontsize{10}{11}\ttfamily,
    %basicstyle=\color{red},
    breaklines=true,%
    morekeywords={matlab2tikz},
    keywordstyle=\color{blue},%
    morekeywords=[2]{1}, keywordstyle=[2]{\color{black}},
    identifierstyle=\color{black},%
    stringstyle=\color{mylilas},
    commentstyle=\color{mygreen},%
    showstringspaces=false,%without this there will be a symbol in the places where there is a space
    numbers=left,%
    numberstyle={\tiny \color{black}},% size of the numbers
    numbersep=9pt, % this defines how far the numbers are from the text
    emph=[1]{for,end,break},emphstyle=[1]\color{red}, %some words to emphasise
    %emph=[2]{word1,word2}, emphstyle=[2]{style},    
}

%--------------------------------------------------------------------------------------
% Main variables
%--------------------------------------------------------------------------------------
\newcommand{\vikszerzoone}{Horny�k M�t� D�niel}
\newcommand{\vikszerzotwo}{Moln�r Marcell}
\newcommand{\vikszerzothr}{Poleczki �kos}
\newcommand{\vikcim}{H�zi feladat tervdokument�ci�}
\newcommand{\viktanszek}{M�r�stechnika �s Inform�ci�s Rendszerek Tansz�k}
\newcommand{\vikdoktipus}{Be�gyazott rendszerek szoftvertechnol�gi�ja}
%\newcommand{\vikdepartmentr}{Moln�r Marcell}

%--------------------------------------------------------------------------------------
% Page layout setup
%--------------------------------------------------------------------------------------
% we need to redefine the pagestyle plain
% another possibility is to use the body of this command without \fancypagestyle
% and use \pagestyle{fancy} but in that case the special pages
% (like the ToC, the References, and the Chapter pages)remain in plane style

\pagestyle{plain}
%\setlength{\parindent}{0pt} % �ttekinthet�bb, angol nyelv� dokumentumokban jellemz�
%\setlength{\parskip}{8pt plus 3pt minus 3pt} % �ttekinthet�bb, angol nyelv� dokumentumokban jellemz�
\setlength{\parindent}{12pt} % magyar nyelv� dokumentumokban jellemz�
\setlength{\parskip}{0pt}    % magyar nyelv� dokumentumokban jellemz�

\chapterfont{\large\upshape\bfseries}
\marginsize{35mm}{25mm}{15mm}{15mm} % anysize package
\setcounter{secnumdepth}{0}
\sectionfont{\large\upshape\bfseries}
\setcounter{secnumdepth}{2}
\singlespacing
\frenchspacing

%--------------------------------------------------------------------------------------
%	Setup hyperref package
%--------------------------------------------------------------------------------------
\hypersetup{
    bookmarks=true,            % show bookmarks bar?
    unicode=true,             % non-Latin characters in Acrobat�s bookmarks
    pdftitle={\vikcim},        % title
    pdfauthor={\vikszerzoone, \vikszerzotwo, \vikszerzothr},    % author
    pdfsubject={\vikdoktipus}, % subject of the document
    pdfcreator={\vikszerzoone, \vikszerzotwo, \vikszerzothr},   % creator of the document
    pdfproducer={Producer},    % producer of the document
    pdfkeywords={keywords},    % list of keywords
    pdfnewwindow=true,         % links in new window
    colorlinks=true,           % false: boxed links; true: colored links
    linkcolor=black,           % color of internal links
    citecolor=black,           % color of links to bibliography
    filecolor=black,           % color of file links
    urlcolor=black             % color of external links
}

%--------------------------------------------------------------------------------------
% Set up listings
%--------------------------------------------------------------------------------------
\lstset{
	basicstyle=\scriptsize\ttfamily, % print whole listing small
	keywordstyle=\color{black}\bfseries\underbar, % underlined bold black keywords
	identifierstyle=, 					% nothing happens
	commentstyle=\color{white}, % white comments
	stringstyle=\scriptsize\sffamily, 			% typewriter type for strings
	showstringspaces=false,     % no special string spaces
	aboveskip=3pt,
	belowskip=3pt,
	columns=fixed,
	backgroundcolor=\color{lightgray},
} 		
\def\lstlistingname{lista}	

%--------------------------------------------------------------------------------------
%	Some new commands and declarations
%--------------------------------------------------------------------------------------
\newcommand{\code}[1]{{\upshape\ttfamily\scriptsize\indent #1}}

% define references
\newcommand{\figref}[1]{\ref{fig:#1}.}
\newcommand{\equatref}[1]{(\ref{eq:#1})}
\newcommand{\listref}[1]{\ref{listing:#1}.}
\newcommand{\sectref}[1]{\ref{sect:#1}}
\newcommand{\tabref}[1]{\ref{tab:#1}.}

\DeclareMathOperator*{\argmax}{arg\,max}
%\DeclareMathOperator*[1]{\floor}{arg\,max}
\DeclareMathOperator{\sign}{sgn}
\DeclareMathOperator{\rot}{rot}
\definecolor{lightgray}{rgb}{0.95,0.95,0.95}

\author{\vikszerzo}
\title{\viktitle}
\includeonly{
	fedlap,%
	chapter1,%
}
%--------------------------------------------------------------------------------------
%	Setup captions
%--------------------------------------------------------------------------------------
\captionsetup[figure]{
%labelsep=none,
font={normalsize,it},
%justification=justified,
width=.75\textwidth,
aboveskip=10pt}

\renewcommand{\captionlabelfont}{\small\bf}
\renewcommand{\captionfont}{\footnotesize\it}

%--------------------------------------------------------------------------------------
% Page numbering for landscape pages
%--------------------------------------------------------------------------------------

\newcommand{\Lpagenumber}{\ifdim\textwidth=\linewidth\else\bgroup
  \dimendef\margin=0 %use \margin instead of \dimen0
  \ifodd\value{page}\margin=\oddsidemargin
  \else\margin=\evensidemargin
  \fi
  \raisebox{\dimexpr -\topmargin-\headheight-\headsep-0.5\linewidth}[0pt][0pt]{%
    \rlap{\hspace{\dimexpr \margin+\textheight+\footskip}%
    \llap{\rotatebox{90}{\thepage}}}}%
\egroup\fi}
\AddEverypageHook{\Lpagenumber}%
%--------------------------------------------------------------------------------------


%--------------------------------------------------------------------------------------
% Table of contents and the main text
%--------------------------------------------------------------------------------------
\begin{document}
\singlespacing

\pagenumbering{arabic}
\onehalfspacing

%--------------------------------------------------------------------------------------
%	The title page
%--------------------------------------------------------------------------------------
\begin{titlepage}
\begin{center}
\includegraphics[width=60mm,keepaspectratio]{figures/BMElogo.png}\\
\vspace{0.3cm}
\textbf{Budapesti M�szaki �s Gazdas�gtudom�nyi Egyetem}\\
\textmd{Villamosm�rn�ki �s Informatikai Kar}\\
\textmd{\viktanszek}\\[5cm]

\vspace{0.4cm}
{\huge \bfseries \vikcim}\\[0.8cm]
\vspace{0.4cm}
{\huge \bfseries \vikcimm}\\[0.8cm]
\vspace{0.5cm}
\textsc{\Large \vikdoktipus}\\[4cm]

\emph{K�sz�tett�k:}\\
\textmd{\vikszerzoone}\\
\textmd{\vikszerzotwo}\\
\textmd{\vikszerzothr}\\

\vfill
{\large \today}
\end{center}
\end{titlepage}


\vfill
%----------------------------------------------------------------------------
\section{Az elk�sz�tend� alkalmaz�s}
%----------------------------------------------------------------------------
A h�zi feladatot be�gyazott k�rnyezetben sz�ks�ges megval�s�tani, ez�rt az Android platformra esett a v�laszt�sunk. Egy k�tj�t�kos, oldaln�zetes, l�v�ld�z�s j�t�k megval�s�t�sa a c�l, ahol egyszerre kett� j�t�kos tud egym�s ellen j�tszani, egy k�z�s lok�lis h�l�zaton (LAN-on).

%----------------------------------------------------------------------------
\section{A j�t�k fel�lete}
%----------------------------------------------------------------------------
Oldaln�zetes, fix kamera�ll�s�, statikus h�tt�rrel rendelkez� fel�let adja az alkalmaz�s kin�zet�t. A saj�t karakter�nkkel val� interakci�t a fel�let jobb �s bal als� sark�ban elhelyezett gombok biztos�tj�k. Ezek k�z�l egy a l�v�s�rt, valamint tov�bbi h�rom a k�t oldalir�ny� mozg�s�rt �s az ugr�s�rt felel�s. Az ir�ny�t�st �t�ll�thatjuk joystickkal t�rt�n� vez�rl�sre is.\\
A p�ly�n megtal�lhat� m�g a kett� j�t�kos karaktere, melyek egy�rtelm�en megk�l�nb�ztethet�k. A karakterek �n. platformokra is felugorhatnak, �gy megnehez�tve a m�sik karakter dolg�t a pontszerz�sben. A mozg�st azonban k�l�nb�z� akad�lyok is nehez�tik, amelyek t�l magasak ahhoz, hogy a karakter �tugorja. Ezeket egy platformra val� r�ugr�ssal, majd az akad�lyra vagy akad�lyon val� �tugr�ssal lehet kiker�lni.

%----------------------------------------------------------------------------
\section{A j�t�k c�lja}
%----------------------------------------------------------------------------
A j�t�kosoknak bizonyos mennyis�g� �lete van. A m�sik j�t�kos l�v�sekkel cs�kkenteni tudja. A l�v�sek csak bizonyos id�k�z�nk�nt enged�lyezettek, melyet a l�v�s gomb letilt�sa val�s�t meg. A l�v�seket egy l�ved�k reprezent�lja a kijelz�n, melynek sebess�ge konstans, �s nagyobb, mint a karakterek maxim�lis sebess�ge.\\
Mindegyik j�t�kos c�lja, hogy min�l t�bbsz�r kinull�zza a m�sik karakter �letpontj�t. Egy kinull�z�s sz�m�t egy pont megszerz�s�nek. A pontszerz�seket sz�moljuk �s aki el�bb �r el egy bizonyos sz�mot az gy�z. A j�t�k elej�n �s minden pontszerz�s ut�n a k�t karakter egy kiindul�si helyre ker�l vissza.\\
Ha a k�t j�t�kos egyszerre null�zza ki a m�sik �let�t, abban az esetben mindegyik egy-egy pontot kap.

%----------------------------------------------------------------------------
\section{A j�t�kosok �sszerendel�se}
%----------------------------------------------------------------------------
Az alkalmaz�s a bel�p�st k�vet�en egy szerverhez kapcsol�dik, amely k�t egym�s ut�n bejelentkezett felhaszn�l�t p�ros�t �ssze. Egy menet addig tart, am�g valaki nem gy�z, vagy ki nem l�p a j�t�kb�l. Minden menet v�g�n �j ellenfelet sorsolunk.

%----------------------------------------------------------------------------
\section{Kezel�fel�let �s egy�b j�t�k be�ll�t�sok}
%----------------------------------------------------------------------------
A m�r kor�bban eml�tett karakter ir�ny�t�s t�rt�nhet joystickkal vagy gombokkal. E k�t opci� k�z�tt a j�t�kosok egy be�ll�t�sok men�ben tudnak v�lasztani. A v�laszt�st az alkalmaz�s megjegyzi, �gy a k�vetkez� bel�p�skor a kor�bban be�ll�tott ir�ny�t�s fog megmaradni. A be�ll�t�sok men�ben a hanger�t szab�lyozhatjuk �s h�tt�rk�pek k�z�tt v�logathatunk m�g.

\label{page:last}
\end{document}
